\documentclass[11pt,a4paper]{article}
\usepackage[T1]{fontenc}
\usepackage[utf8]{inputenc}
\usepackage{amsmath,amssymb,amsthm}
\usepackage[margin=1in]{geometry}
\usepackage[colorlinks=true,linkcolor=blue,urlcolor=blue]{hyperref}
\theoremstyle{plain}
\newtheorem{theorem}{Theorem}
\newtheorem{lemma}[theorem]{Lemma}
\newtheorem{corollary}[theorem]{Corollary}
\theoremstyle{definition}
\newtheorem{remark}[theorem]{Remark}
\title{A proof of the conjectured recurrence for OEIS A348410}
\author{Adrian Perez Fontelles\\ \small Independent researcher}
\date{30 August 2026}
\begin{document}
\maketitle

\begin{abstract}
OEIS A348410 carries a conjectured linear recurrence with polynomial coefficients, of
order 4. The entry posts no generating function; it defines the sequence by an
extraction, $a(n)=[x^{n}]\,f(x)g(x)^{n}$. We derive a generating function from that
definition and prove the conjecture from it. Reading the extraction as a contour integral
and summing the geometric series turns $\sum_{n}a(n)t^{n}$ into a single integral whose
only pole inside a small circle is the branch of $x=t\,g(x)$ through the origin; the
residue there gives $A(t)=f(x(t))/\bigl(1-t\,g'(x(t))\bigr)$. So $A$ is algebraic, with
minimal polynomial obtained by eliminating $x$, and the recurrence follows from the
residual test in the field that polynomial generates. Here the residual is
$2010 x^{4} + 47040 x^{3} - 38400 x^{2}$, of degree $4$, which proves the recurrence for every $n>4$.
\end{abstract}

\noindent\small 2020 Mathematics Subject Classification. 05A15, 11B37, 30B10, 33F10.\normalsize

\section{The sequence and the conjecture}

OEIS A348410 is ``Number of nonnegative integer solutions to n = Sum\_\{i=1..n\} (a\_i + b\_i), with b\_i even''. It has offset $0$ and begins
\[
a(0),\dots,a(7)\;=\;1, 1, 5, 19, 85, 376, 1715, 7890,\ \dots
\]
The entry states, not as a conjecture,
\begin{quote}\small
a(n) = [x\textasciicircum{}n] ( (1 - x)*(1 - x\textasciicircum{}2) )\textasciicircum{}(-n).
\end{quote}
that is,
\begin{equation}\label{eq:def}
a(n)\;=\;[x^{n}]\,f(x)\,g(x)^{n},
\qquad f(x)=1,\qquad g(x)=\frac{1}{x^{3} - x^{2} - x + 1} .
\end{equation}
This is the only input taken from the entry, and it was checked against the entry's own
DATA before use. In particular \emph{no generating function is assumed}: the entry posts
none, and the one used below is derived.

The entry separately carries the following comment:
\begin{quote}\small
\textbf{Conjecture:} D-finite with recurrence +7168*n*(2*n-1)*(n-1)*a(n) -64*(n-1)*(1759*n\textasciicircum{}2-5294*n+5112)*a(n-1) +12*(7561*n\textasciicircum{}3-75690*n\textasciicircum{}2+165271*n-101070)*a(n-2) +5*(110593*n\textasciicircum{}3-743946*n\textasciicircum{}2+1659971*n-1232778)*a(n-3) +2680*(4*n-15)*(2*n-7)*(4*n-13)*a(n-4)=0. - \emph{R. J. Mathar}, Oct 19 2021
\end{quote}
As of the ``Last modified'' line on the live entry (2026-06-29, revision 60) the
statement is still recorded as a conjecture, and no proof appears among the entry's links,
formulas or programs.

\section{A generating function from the extraction}

\begin{lemma}\label{lem:diag}
Let $f,g$ be rational functions regular at $0$ with $g(0)\neq0$, and let
$a(n)=[x^{n}]f(x)g(x)^{n}$. Then $A(t)=\sum_{n\ge0}a(n)t^{n}$ satisfies
\[
A(t)\;=\;\frac{f\bigl(x(t)\bigr)}{1-t\,g'\bigl(x(t)\bigr)},
\]
where $x(t)$ is the unique power series with $x(0)=0$ solving $x=t\,g(x)$. In particular
$A$ is algebraic over $\mathbb{Q}(t)$.
\end{lemma}

\begin{proof}
Write the extraction as a contour integral on a circle $|x|=\rho$ small enough that $f$
and $g$ are regular on and inside it:
\[
a(n)=\frac{1}{2\pi i}\oint \frac{f(x)g(x)^{n}}{x^{n+1}}\,dx .
\]
For $|t|$ small, $|t\,g(x)/x|<1$ on the circle, so the geometric series may be summed
under the integral:
\[
A(t)=\frac{1}{2\pi i}\oint\frac{f(x)}{x}\sum_{n\ge0}\Bigl(\frac{t\,g(x)}{x}\Bigr)^{n}dx
     =\frac{1}{2\pi i}\oint\frac{f(x)}{x-t\,g(x)}\,dx .
\]
Since $g(0)\neq0$, the equation $x=t\,g(x)$ has, for small $t$, exactly one root in the
disc, namely the branch $x(t)$ with $x(0)=0$ furnished by Lagrange inversion; it is a
simple zero of $x-t\,g(x)$ because the derivative $1-t\,g'(x)$ is $1$ at $t=0$. The
residue theorem gives the stated value. Both $x(t)$ and hence $A(t)$ are algebraic,
being defined by polynomial equations over $\mathbb{Q}(t)$.
\end{proof}

\begin{corollary}\label{cor:minpoly}
$A$ satisfies the polynomial obtained by eliminating $x$ between
$x-t\,g(x)=0$ and $y\bigl(1-t\,g'(x)\bigr)-f(x)=0$.
\end{corollary}

\begin{proof}
Both equations hold at $x=x(t)$, $y=A(t)$ by Lemma~\ref{lem:diag}. Eliminating $x$ by a
resultant produces a polynomial in $t$ and $y$ vanishing there.
\end{proof}

For A348410 this yields
\[
256 t^{2} y^{4} - 256 t^{2} y^{3} + 96 t^{2} y^{2} - 16 t^{2} y + t^{2} + 107 t y^{4} - 107 t y^{3} + 36 t y^{2} - 4 t y - 32 y^{4} + 32 y^{3} \;=\;0 ,
\]
and $A$ is the branch of this polynomial whose expansion is the entry's own sequence.
The branch is identified by that expansion, not by solving for it in radicals --- the
polynomial can have degree beyond four, where no such expression exists.

\section{The recurrence}

Let $\theta=x\,\dfrac{d}{dx}$, acting on $x^{m}$ as multiplication by $m$.

\begin{lemma}\label{lem:transfer}
Let $p_{0},\dots,p_{r}$ be polynomials and put $b(n)=\sum_{i}p_{i}(n)a(n-i)$, with
$a(m)=0$ for $m<0$. Then $B(x)=\sum_{n}b(n)x^{n}=\sum_{i}x^{i}(p_{i}(\theta+i)A)(x)$, so
the recurrence $\sum_{i}p_{i}(n)a(n-i)=0$ holds for every $n>d$ if and only if $B$ is a
polynomial of degree at most $d$.
\end{lemma}

\begin{proof}
Substituting $m=n-i$ gives $\sum_{n}p_{i}(n)a(n-i)x^{n}=x^{i}(p_{i}(\theta+i)A)(x)$, and
$b(n)$ is the coefficient of $x^{n}$ in $B$.
\end{proof}

Since $A$ lies in the algebraic function field generated by the polynomial of
Corollary~\ref{cor:minpoly}, and such a field is closed under $\theta$ --- differentiating
its defining relation expresses the derivative of the generator inside the field ---
every $p_{i}(\theta+i)A$ lies there too, and so does $B$. Computing it exactly gives
\[
B(x)\;=\;2010 x^{4} + 47040 x^{3} - 38400 x^{2} ,
\]
a polynomial of degree $4$.

\begin{theorem}
The conjectured recurrence
\[
7168 n \left(n - 1\right) \left(2 n - 1\right)\,a(n) -  64 \left(n - 1\right) \left(1759 n^{2} - 5294 n + 5112\right)\,a(n-1) + 12 \left(7561 n^{3} - 75690 n^{2} + 165271 n - 101070\right)\,a(n-2) + 5 \left(110593 n^{3} - 743946 n^{2} + 1659971 n - 1232778\right)\,a(n-3) + 2680 \left(2 n - 7\right) \left(4 n - 15\right) \left(4 n - 13\right)\,a(n-4) \;=\;0
\]
holds for every $n>4$.
\end{theorem}

\begin{proof}
Immediate from Lemma~\ref{lem:transfer} and the displayed value of $B$.
\end{proof}

\section{Verification}

Every number above is an output of a script written separately from this note, and three
independent checks were run.

First, the derived generating function was not assumed to be right. Its series was
computed from \eqref{eq:def} in two independent ways --- once through the residue formula
of Lemma~\ref{lem:diag}, by solving $x=t\,g(x)$ as a truncated series, and once directly,
by expanding $f\,g^{n}$ and reading off the coefficient of $x^{n}$ --- and the two agree.
The result was then compared term by term with the DATA section of A348410, agreeing on
all 13 terms compared.

Second, the recurrence was evaluated directly on the published terms, in exact integer
arithmetic and without reference to any generating function: for each $n$ in range the sum
$\sum_{i}p_{i}(n)a(n-i)$ was formed from the entry's own values and found to vanish, for
all 21 values of $n$ from $n=5$ upward.

Third, the symbolic step is exact throughout. The residual is obtained by
rational-function arithmetic in the algebraic function field, not by series truncation, so
the identity $B=2010 x^{4} + 47040 x^{3} - 38400 x^{2}$ is an identity of functions and the theorem holds for all
$n>4$ simultaneously.

\begin{thebibliography}{9}
\bibitem{oeis} The OEIS Foundation, \emph{The On-Line Encyclopedia of Integer Sequences},
\url{https://oeis.org}, 2026. Sequence A348410.
\bibitem{stanley} R.~P.~Stanley, \emph{Enumerative Combinatorics}, Volume 2, Cambridge
University Press, 1999, Section 6.3 (Lagrange inversion) and Chapter 6 (algebraic and
D-finite generating functions).
\bibitem{flajolet} P.~Flajolet and R.~Sedgewick, \emph{Analytic Combinatorics}, Cambridge
University Press, 2009, Chapter VII.
\bibitem{kp} M.~Kauers and P.~Paule, \emph{The Concrete Tetrahedron}, Springer, 2011,
Chapter 7.
\end{thebibliography}
\end{document}
